\subsection{Tauschermodule gebaut und mit FluidX3D durchströmt}\label{sec:tauschercfd}
Die Bauartwahl des Vorwärmers stützt sich bisher auf Korrelationen aus der Literatur, deren Gültigkeitsbereich
($\Pr = 3$--5) das Divertorwasser mit $\Pr \approx 1{,}3$--1,8 nicht trifft. Deshalb wurden die Module gebaut und
die Strömung selbst gerechnet. OpenFOAM schafft der vorhandene Rechner nicht; gerechnet wird mit dem
Lattice-Boltzmann-Löser FluidX3D auf der Grafikkarte \cite{fluidx3d,lehmann2022fp}.

\paragraph{1. Module.} PicoGK baut je einen abgestuften Zwei-Fluid-Block ($48\times48\times96$\,mm, Wand \SI{0,5}{mm},
Hülle \SI{1,5}{mm}): Gyroid mit Zellgröße \SI{4,33}{mm} $\to$ \SI{5,8}{mm} und Diamond mit \SI{5,37}{mm} $\to$
\SI{7,2}{mm}. Beide ergeben denselben hydraulischen Durchmesser $d_h = 2 \to 2{,}9$\,mm. Die Zellgröße wächst
logarithmisch längs des Blocks, die Phase dazu ist $p_z = 2\pi L/(a_1 - a_0)\,\ln(a/a_0)$. Endkappen schließen die
kalte Seite, Seitenfenster die warme. Bei \SI{0,2}{mm} Voxelgröße sind 39,7\,\% Stahl, je 30\,\% warme und kalte
Seite. Die Platten (glatt und Chevron, $\beta = \pm 60^\circ$) entstehen parametrisch in FreeCAD.

\bild{tauscher_module.jpg}{1.0}{Abgestufte Zwei-Fluid-Module, aufgeschnitten. Links Gyroid, rechts Diamond. Rot die
warme (Divertor-), blau die kalte (Kondensat-)Seite, grau die Stahlwand.}
\lesehilfe{Zwei aufgeschnittene Tauscherblöcke, links mit Gyroid-, rechts mit Diamond-Gitter (zwei
verschiedene Gitterformen). Rot ist der Raum für das warme Wasser aus dem Divertor, blau der für das kalte
Kondensat, grau die Stahlwand dazwischen. Es sind zwei getrennte, ineinander verschlungene Kanalsysteme; die beiden
Wässer berühren sich nie, nur die Wärme geht durch die Wand.}


\paragraph{2. Rechengebiet.} Ein Gebiet mit Ein- und Auslass scheitert bei $\mathrm{Re}\approx3000$: Der nötige
Druckabfall verlangt im Lattice-Boltzmann-Verfahren mehr als 25\,\% Dichteunterschied, weit außerhalb der
Gültigkeit. Gerechnet wird deshalb ein in allen Richtungen periodischer Ausschnitt einer Seite. Eine Volumenkraft
$f_z$ treibt die Strömung, eine gleichmäßige Wärmequelle $q$ heizt das Fluid, die Wand hält $T_w=0$. Daraus folgen
\[
f = \frac{f_z\,d_h}{2\,u^2},\qquad h = \frac{q\,V_f}{A\,T_m},\qquad \mathrm{Nu} = \frac{h\,d_h}{D},\qquad
T_m = \frac{\sum u_z T}{\sum u_z}.
\]
$f_z$ wird iteriert, bis Re auf 15\,\% stimmt.

\paragraph{3. Änderungen an FluidX3D.} Fünf Punkte mussten geändert werden (Quellen mit Lizenzhinweis in
\texttt{fluidx3d/}):
\begin{itemize}
\item Das Temperaturmodell (D3Q7) hat $c_s^2 = 1/4$. Die tatsächliche Diffusivität ist $D = \alpha/2$; ein reiner
  Diffusionstest bestätigte das (gemessen $D/\alpha = 0{,}44$--0,53). Übergeben wird daher $\alpha = 2D$.
\item Die Wand lag eine Zelle zu weit im Fluid, wenn die Wandtemperatur in Fluidzellen gesetzt wurde. Das ergab ein
  etwa dreimal zu hohes Nu. Jetzt gilt an Festzellen ein Anti-Bounce-Back: $g_\text{ein} = -g_\text{aus} +
  \tfrac14(T_w - 1)$.
\item Bei $\mathrm{Re}=3000$ liegt die Zell-Péclet-Zahl bei etwa 185. Aus dem Smagorinsky-Modell des Lösers kommt
  deshalb eine turbulente Diffusivität $D_t = \nu_t/0{,}85$ dazu.
\item Der einfache BGK-Stoß der Temperatur wurde bei 96 Zellen je $d_h$ nach \num{172000} Schritten schlagartig
  instabil (Temperatur in 500 Schritten auf $10^{20}$). Die Relaxationszeit liegt mit $\tau \approx 0{,}504$ dicht an
  der Stabilitätsgrenze 0,5. Ersetzt wurde er durch einen TRT-Stoß nach Ginzburg \cite{ginzburg2005}: Der antisymmetrische Anteil
  relaxiert mit der Rate, die $D$ festlegt, der symmetrische mit der „magischen“ Kombination
  $\Lambda = (1/\omega^+ - \tfrac12)(1/\omega^- - \tfrac12) = 1/4$.
\item Die Wärmequelle ist wahlweise gleichmäßig oder proportional zu $u_z$. Die zweite Form ist die übliche
  Formulierung für voll ausgebildete periodische Strömung mit axial gleichmäßigem Wärmestrom (Patankar, Liu und
  Sparrow 1977 \cite{patankar1977}). Eine gleichmäßige Quelle erhöht Nu schon im laminaren Spalt von 8,235 auf 10.
\end{itemize}

\paragraph{4. Prüfung.} Für den laminaren Plattenkanal sind die Werte exakt bekannt: $f\,\mathrm{Re}=24$, dazu
$\mathrm{Nu}=10$ (gleichmäßige Quelle) bzw.\ 8,235 (Quelle $\propto u$). Die Rechnung liefert $f\,\mathrm{Re} = 29{,}7
\to 25{,}7 \to 24{,}7$ bei 12, 24, 48 Zellen je Kanalhöhe. Für Nu ergeben sich $10{,}6 \to 10{,}3$ (BGK, gleichmäßig)
und $8{,}9 \to 8{,}5$ (TRT, $\propto u$). Der Fehler fällt mit der Ordnung $p \approx 1{,}3$. Die Wahl von $\Lambda$
(1/4 oder 3/16) änderte den Prüffall nicht.

\bild{tauscher_cfd_pruefung.png}{1.0}{Links: Prüffall laminarer Plattenkanal, Abweichung vom exakten Wert. Rechts:
Nusselt-Zahl bezogen auf Renon bei derselben Re, über der Gitterfeinheit; offene Kreise: Richardson-Schätzung.}
\lesehilfe{Links eine Prüfung des Strömungsprogramms an einem Fall mit exakt bekannter
Lösung: Strömung zwischen zwei Platten. Senkrecht steht die Abweichung in Prozent, waagrecht, wie fein das
Rechengitter ist. Je feiner, desto genauer; bei 48 Zellen sind es etwa 3\,\%. Rechts die \emph{Nusselt-Zahl},
ein Maß dafür, wie gut Wärme von der Wand ins Wasser übergeht, geteilt durch den Wert aus der Literatur (Renon, $=1$).
Die offenen Kreise rechnen auf ein unendlich feines Gitter hoch (Richardson) und liegen etwa 25\,\% über der
Literatur. Das graue Kreuz stammt aus einem älteren, zu ungenauen Verfahren.}


\paragraph{5. Vergleich der Bauarten.} Alle Rechnungen mit TRT und Quelle $\propto u_z$, $\Pr = 1{,}5$. Der Antrieb
wird nachgeregelt, bis Re auf 15\,\% stimmt; bei 96 Zellen blieben Gyroid und Diamond etwas darunter. $j/f^{1/3}$
mit $j = \mathrm{Nu}/(\mathrm{Re}\,\Pr^{1/3})$ ist die Flächengüte: wie viel Wärmeübergang je Pumpaufwand.

\begin{table}[H]\centering\small
\begin{tabular}{llrrrrr}\toprule
Bauart & Gitter & Re & $f$ & Nu & $j/f^{1/3}$ & Nu/Nu$_\text{Renon}$ \\\midrule
Gyroid  & 48 & 2952 & 0,805 & 139,8 & 0,0445 & 1,86 \\
Gyroid  & 96 & 2583 & 0,732 & 102,0 & 0,0383 & 1,49 \\
Diamond & 48 & 2974 & 0,668 & 133,4 & 0,0448 & 1,77 \\
Diamond & 96 & 2676 & 0,596 & 100,9 & 0,0391 & 1,44 \\
Chevron $\pm60^\circ$ & 48 & 2908 & 1,006 & 121,5 & 0,0364 & -- \\
Chevron $\pm60^\circ$ & 96 & 2921 & 0,863 & 103,5 & 0,0325 & -- \\\bottomrule
\end{tabular}
\caption{FluidX3D-Ergebnisse je Seite (periodischer Ausschnitt). Renon gilt nur für TPMS.}
\end{table}

\paragraph{Befund.}
\begin{itemize}
\item \textbf{Hydraulik:} $f$ liegt im Band der Korrelationen (Renon 0,48--0,56, Li 2022 0,81--0,83
  \cite{renon2024,alsafadi2026}). Diamond hat bei gleichem $d_h$ etwa 20\,\% weniger Reibung als Gyroid; Renon gibt
  12--17\,\%.
\item \textbf{Wärmeübergang:} Nu liegt über den Korrelationen, aber der Abstand schrumpft mit dem Gitter: Faktor 1,8
  bei 48 und 1,5 bei 96 Zellen. Mit der Ordnung aus dem Prüffall ergibt eine Richardson-Schätzung etwa 1,2--1,25. Das
  war eine grobe Schätzung aus zwei Stufen. Die dritte Stufe (144 Zellen, 62,5 Mio.\ Zellen) ergab 1,36 (Gyroid) und
  1,30 (Diamond) und zeigt, dass die Reihe noch nicht asymptotisch ist; extrapoliert liegt Nu bei Renon
  (Abschnitt~\ref{sec:tauschersurrogat}).
\item \textbf{Stoßmodell:} BGK hatte Nu auch ohne Absturz um etwa 25\,\% überschätzt (194 statt 154 bei 48 Zellen,
  gleiche Quelle).
\item \textbf{Bauarten:} Gyroid und Diamond sind bei der Flächengüte gleichauf (0,038 bzw.\ 0,039 bei 96 Zellen).
  Chevron-Platten liegen etwa 15\,\% darunter: Sie erreichen ähnliches Nu nur mit mehr Reibung. Das stützt die
  Rangfolge aus Abschnitt~\ref{sec:vorwaermer} (TPMS kompakter als Platten). Einen Vorteil des Gyroids gegenüber
  Diamond zeigt die Rechnung nicht.
\item \textbf{Grenzen:} gleichmäßig große Zellen (die Stufung wurde nicht gerechnet), nur eine Re, $\Pr = 1{,}5$,
  Wand isotherm und ohne Wärmeleitung im Stahl, Turbulenz über ein LES-Modell bei nur mäßiger Auflösung.
\end{itemize}

\bild{tauscher_stroemung_drei.jpg}{1.0}{Teilchenbahnen in Gyroid, Diamond und Chevron-Platten (von links), Re 3000,
48 Zellen je $d_h$. Die Wand ist im Röntgenblick dargestellt. Aus den Bahnen entstand je ein Video. Die Farbe der
Bahnen stammt aus den Läufen mit BGK-Stoß und dient nur der Anschauung.}
\lesehilfe{Standbilder der Bahnen von Wasserteilchen (helle Linien) durch drei Bauarten:
Gyroid, Diamond und Chevron-Platten. Die Wand ist durchscheinend dargestellt. Je stärker sich die Bahnen verwirbeln
und kreuzen, desto besser mischt sich das Wasser und desto mehr Wärme nimmt es von der Wand auf. Die Farbe hat
hier keine Bedeutung.}


Aus diesen und weiteren 30 Läufen entstand das Ersatzmodell in Abschnitt~\ref{sec:tauschersurrogat}.
